% Time form: Past tense

\newpage
\section{Self-Hosted \glsxtrshort{LLM}}\label{sec:self_hosted_llm}
During the risk assessment~\ref{subsec:risk_assessment} we identified that the \glsxtrshort{AI} tokens required to use cloud based \glsxtrshort{LLM}s might be very expensive, and it's really hard to estimate the costs.
One of the possibilities to avoid this risk was the usage of self-hosted \glsxtrshort{LLM}s, and luckily we already had existing hardware at home which is suitable for running smaller \glsxtrshort{LLM}s.
Besides eliminating the costs for \glsxtrshort{AI} tokens, at least during the development phase, it also allows the quick evaluation of different \glsxtrshort{AI} Models.
Depending on the quality and performance of the smaller models running in the self-hosted setup we can consider not using a cloud based \glsxtrshort{LLM} at all.

\subsection{Hardware Setup}\label{subsec:hw_setup}

\begin{table}[h]
    \centering
    \begin{tabular}{ll}
        \hline
        \textbf{Component} & \textbf{Specification} \\
        \hline
        CPU & AMD Ryzen 5\ 3600 (6 cores / 12 threads, up to 4.2 GHz) \\
        GPU & NVIDIA GeForce RTX 3070 (8 GB \glsxshort{VRAM}) \\
        System Memory & 32 GB RAM \\
        \hline
    \end{tabular}
    \caption{Hardware configuration of the self-hosted \glsxtrshort{LLM} inference server.}
    \label{tab:llm_hardware_setup}
\end{table}

The 8 GB \glsxshort{VRAM} of the RTX 3070 allows us to run models with up to 7 billion parameters, depending on the specific model~\cite{simplr2025vramllm}.

\subsection{System Setup}\label{subsec:system_setup}


\begin{lstlisting}[language=bash,caption={UFW configuration},label={lst:llm-base-system-setup}]
# Create a new user for isolating the services
sudo useradd -m ai

# Enable linger for user ai, required to run services such as rootles containers without being logged in
sudo loginctl enable-linger ai

systemctl --user daemon-reexec

# Install podman
sudo pacman -Syu podman

# Login as user ai
sudo su ai

\end{lstlisting}

\subsubsection{LLM Model Manager - Ollama}\label{subsubsec:llm_model_manager_ollama}

\begin{lstlisting}[language=bash,caption={UFW configuration},label={lst:ollama-setup}]
# Install Ollama

\end{lstlisting}

\subsubsection{Webserver - Nginx}\label{subsubsec:webserver_nginx}



\subsubsection{LLM Chat GUI - Open WebUI}\label{subsubsec:open_webui}



\subsubsection{System Hardening}\label{subsubsec:system_hardening}


\begin{lstlisting}[language=bash,caption={UFW configuration},label={lst:ufw-configuration}]
# Install UFW
sudo pacman -Syu ufw
sudo ufw status

# Allow all outgoing connections
sudo ufw default allow outgoing

# Deny all incoming connections per default
sudo ufw default deny incoming

# Allow SSH for remote access
sudo ufw allow ssh

# Allow HTTP for certbot HTTP challange
sudo ufw allow http

# Allow HTTPS for Open WebUI
sudo ufw allow https

# Allow default Ollama port
sudo ufw allow 11434/tcp

# Make sure to run this command after adding the "allow" rules!
sudo ufw enable

sudo ufw status
\end{lstlisting}

\subsubsection{}\label{subsubsec:}
Port Forwarding via Fritzbox
DynDNS via Infomaniak

\subsection{Evaluation}\label{subsec:evaluation}
